\documentclass{article}
\usepackage[top=2.3cm, bottom=2.3cm, left=2.1cm, right=2.1cm]{geometry}
\usepackage{amsmath,amssymb,amsthm,mathtools}
\usepackage{bbm}
\usepackage{hyperref}
\usepackage[default,regular,black]{sourceserifpro}
\usepackage[T1]{fontenc}
\hypersetup{colorlinks,linkcolor={blue},citecolor={blue},urlcolor={red}}

\newcommand{\Align}{\mathrm{Align}}
\newcommand{\FA}{\mathrm{FA}}
\newcommand{\ind}{\mathbf{1}}
\newcommand{\cF}{\mathcal{F}}
\newcommand{\cFpre}{\mathcal{F}^{\mathrm{pre}}}

\theoremstyle{plain}
\newtheorem{proposition}{Proposition}
\newtheorem{remark}{Remark}

\begin{document}

\title{Review Note: Aggregation from Firm-Level to Sector-Level\\
Extensive Margin (Levels)}
\author{}
\date{\today}
\maketitle

\section{Setup and Notation}

This note reviews the claim in Section~2.1 (``Extensive margin (levels)'')
that the firm-level and sector-level extensive-margin regressions recover the
same $\lambda^\mathrm{ext}_\ell$ coefficient when every firm receives equal
weight.  We verify the aggregation algebra, identify the exact conditions under
which the claim holds, and clarify what happens when the firm-level regression
uses employment weights instead.

\medskip\noindent\textbf{Firm-level regression.}
For a generic cell $c = (j,m,e)$ containing the pre-election firm set
$\cFpre_{jme}$ with $N_c \equiv |\cFpre_{jme}|$ firms, and for each
post-election year $t \in \{e+1,\dots,e+4\}$, the paper specifies
\begin{equation}\label{eq:firm}
  \ind(\text{BNDES}_{fmt} > 0)
  = \sum_\ell \lambda^\mathrm{ext}_\ell\, \FA^\ell_{fmt}
  + \gamma_f + \alpha_{mt} + u_{fmt}.
\end{equation}

\noindent\textbf{Sector-level aggregation.}
The paper defines
\begin{align}
  H^\mathrm{pre}_{jmt}
  &\equiv \frac{1}{N_c}\sum_{f\in\cFpre_{jme}} \ind(\text{BNDES}_{fmt}>0),
  \label{eq:H}\\[4pt]
  \overline{\FA}^{\ell,\mathrm{pre}}_{jmt}
  &\equiv \frac{1}{N_c}\sum_{f\in\cFpre_{jme}} \FA^\ell_{fmt},
  \label{eq:FAbar}
\end{align}
and claims that averaging~\eqref{eq:firm} across the pre-election firm base
yields
\begin{equation}\label{eq:sector}
  H^\mathrm{pre}_{jmt}
  = \sum_\ell \lambda^\mathrm{ext}_\ell\,
    \overline{\FA}^{\ell,\mathrm{pre}}_{jmt}
  + \bar\gamma_{jme} + \alpha_{mt} + \epsilon_{jmt},
\end{equation}
with
$\bar\gamma_{jme} = N_c^{-1}\sum_{f\in\cFpre_{jme}}\gamma_f$
and
$\epsilon_{jmt} = N_c^{-1}\sum_{f\in\cFpre_{jme}} u_{fmt}$.

\section{Verification of the Averaging Step}

\begin{proposition}[Algebraic correctness of the averaging]
If equation~\eqref{eq:firm} holds for every $f\in\cFpre_{jme}$, then
summing both sides over $f$ and dividing by $N_c$ yields
equation~\eqref{eq:sector} exactly, with the same $\lambda^\mathrm{ext}_\ell$.
\end{proposition}

\begin{proof}
Sum~\eqref{eq:firm} over $f\in\cFpre_{jme}$:
\[
  \sum_f \ind(\text{BNDES}_{fmt}>0)
  = \sum_\ell \lambda^\mathrm{ext}_\ell \sum_f \FA^\ell_{fmt}
  + \sum_f \gamma_f
  + N_c\,\alpha_{mt}
  + \sum_f u_{fmt}.
\]
Divide both sides by $N_c$:
\[
  \underbrace{\frac{1}{N_c}\sum_f \ind(\text{BNDES}_{fmt}>0)}_{=\, H^\mathrm{pre}_{jmt}}
  = \sum_\ell \lambda^\mathrm{ext}_\ell
    \underbrace{\frac{1}{N_c}\sum_f \FA^\ell_{fmt}}_{=\,\overline{\FA}^{\ell,\mathrm{pre}}_{jmt}}
  + \underbrace{\frac{1}{N_c}\sum_f \gamma_f}_{=\,\bar\gamma_{jme}}
  + \alpha_{mt}
  + \underbrace{\frac{1}{N_c}\sum_f u_{fmt}}_{=\,\epsilon_{jmt}}.
\]
This is exactly~\eqref{eq:sector}.  The $\lambda^\mathrm{ext}_\ell$
coefficients pass through unchanged because averaging is a linear operation
and the coefficients are common to all firms.
\end{proof}

\begin{remark}[The $\alpha_{mt}$ term]
The municipality$\times$year fixed effect $\alpha_{mt}$ does not carry a firm
subscript, so $\frac{1}{N_c}\sum_f \alpha_{mt} = \alpha_{mt}$.  This is
correct in the paper.
\end{remark}

\section{Under What Conditions Are the Estimated Coefficients Identical?}

The algebraic identity above shows that the \emph{population} equation
aggregates exactly.  However, the \emph{estimated} $\hat\lambda$ from OLS on
the firm-level data and from OLS on the cell-averaged data need not be
numerically equal.  The key issue is the weighting implicit in each estimator.

\subsection{Case 1: Unweighted firm-level OLS}

In the unweighted firm-level regression~\eqref{eq:firm}, every
firm-observation $(f,t)$ receives weight~1.  By the Frisch--Waugh--Lovell
(FWL) theorem, after absorbing fixed effects $(\gamma_f, \alpha_{mt})$, OLS
minimizes
\[
  \sum_c \sum_{f\in\cFpre_{jme}} \Bigl(
    \tilde{Y}_{fmt} - \sum_\ell \lambda_\ell\,\tilde{X}^\ell_{fmt}
  \Bigr)^2,
\]
where tildes denote residuals after the FE projection.  Each cell $c$
contributes $N_c$ terms to the sum.

Now consider running OLS on the cell-averaged equation~\eqref{eq:sector},
where each cell $(j,m,t)$ is one observation.  Each cell gets weight~1, so
OLS minimizes
\[
  \sum_c \Bigl(
    \tilde{\bar{Y}}_c - \sum_\ell \lambda_\ell\,\tilde{\bar{X}}^\ell_c
  \Bigr)^2.
\]
Here each cell contributes exactly~1 term regardless of how many firms it
contains.

\begin{proposition}[Equal coefficients iff cell-size weights]\label{prop:weights}
Let $\hat\lambda^\mathrm{firm}$ be the OLS estimator from the unweighted
firm-level regression~\eqref{eq:firm}, and let $\hat\lambda^\mathrm{sector}$
be the OLS estimator from the cell-averaged regression~\eqref{eq:sector}.
Then:
\[
  \hat\lambda^\mathrm{firm} = \hat\lambda^\mathrm{sector}
  \quad\Longleftrightarrow\quad
  \text{the sector-level regression weights each cell by } N_c.
\]
\end{proposition}

\begin{proof}
Consider the firm-level residuals after FE projection.  For firm $f$ in
cell $c$, the within-cell residual is $\tilde{Y}_{fmt} = Y_{fmt} - \bar{Y}_c
- \bar{Y}_f + \overline{\overline{Y}}$, and similarly for $X$.

We can decompose the firm-level sum of squares into a \emph{between-cell}
component and a \emph{within-cell} component.  The between-cell normal equation
from the firm-level OLS is
\[
  \sum_c N_c \,\tilde{\bar{X}}_c'\bigl(\tilde{\bar{Y}}_c
  - \tilde{\bar{X}}_c\,\lambda\bigr) = 0,
\]
i.e., the cell means $\tilde{\bar{X}}_c$, $\tilde{\bar{Y}}_c$ enter with
weight~$N_c$.  An unweighted regression on cell means weights every cell
equally.  The two coincide if and only if $N_c$ is used as an analytic weight
in the sector-level regression.

Intuitively: if cell $A$ has 10 firms and cell $B$ has 1 firm, the firm-level
regression lets cell~$A$ contribute 10 times the moment conditions.
The cell-level regression with unit weights treats them equally.  Only
weighting by $N_c$ restores the equivalence.
\end{proof}

\subsection{Case 2: Employment-weighted firm-level OLS}

The primary specification in the paper uses analytic regression weights
$n_{\mathrm{emp},ft}$.  The firm-level estimator then solves
\[
  \sum_c \sum_{f\in\cFpre_{jme}} n_{\mathrm{emp},ft}\,
  \tilde{X}^\ell_{fmt}\bigl(\tilde{Y}_{fmt}
  - \sum_\ell \lambda_\ell\,\tilde{X}^\ell_{fmt}\bigr) = 0.
\]
In this case, even weighting the sector-level regression by $N_c$ does
\emph{not} replicate the firm-level estimator, because within each cell
employment-weighted averaging gives more influence to larger firms.

To recover the employment-weighted firm-level coefficient at the sector level,
one would need to:
\begin{enumerate}
  \item Define employment-weighted cell averages:
    \[
      H^\mathrm{emp}_{jmt}
      = \frac{\sum_f n_{\mathrm{emp},ft}\,\ind(\text{BNDES}_{fmt}>0)}
             {\sum_f n_{\mathrm{emp},ft}},
      \qquad
      \overline{\FA}^{\ell,\mathrm{emp}}_{jmt}
      = \frac{\sum_f n_{\mathrm{emp},ft}\,\FA^\ell_{fmt}}
             {\sum_f n_{\mathrm{emp},ft}},
    \]
    where sums run over $f\in\cFpre_{jme}$.
  \item Weight the sector-level regression by total cell employment
    $E_c = \sum_{f\in\cFpre_{jme}} n_{\mathrm{emp},ft}$.
\end{enumerate}
Only then would $\hat\lambda^\mathrm{sector} = \hat\lambda^\mathrm{firm}$.

The paper's current sector-level aggregation uses \emph{equal-firm} (simple
average) weights, $1/N_c$, which corresponds to the \emph{unweighted}
firm-level estimator, not the employment-weighted one.  This is not an error
--- it is a different estimand.  The equal-firm sector average asks about the
average firm; the employment-weighted version asks about the average worker's
firm.

\section{The ``Most Cells Have 1--2 Firms'' Argument}

The user raises the following intuition: since most municipality-sector cells
contain only 1 or 2 firms, cells with more firms should receive proportionally
more weight in the sector regression so that each \emph{firm} carries the same
influence.  This is exactly Proposition~\ref{prop:weights}.

\begin{remark}[Practical magnitude of the discrepancy]
If most cells have $N_c = 1$, then $N_c$-weighting barely changes
anything---for singleton cells, the cell average \emph{is} the firm
observation.  The discrepancy comes entirely from the few cells with $N_c > 1$.
In those cells:
\begin{itemize}
  \item The \emph{unweighted} sector regression treats the cell as one
    observation (the average of, say, 5 firms gets weight~1).
  \item The $N_c$-weighted sector regression gives that cell weight~5.
  \item The firm-level regression gives each of the 5 firms weight~1,
    for a total cell contribution of~5.
\end{itemize}
So the $N_c$-weighted sector regression reproduces the firm-level OLS exactly,
while the unweighted version under-represents the larger cells.

Whether the discrepancy matters empirically depends on:
\begin{enumerate}
  \item The fraction of observations in multi-firm cells.
  \item Whether the within-cell variation in $\FA^\ell_{fmt}$ is large
    relative to between-cell variation.
\end{enumerate}
If most identifying variation is between cells rather than within cells, the
two estimators will be close even without $N_c$-weighting.
\end{remark}

\section{Summary of the Claim in Line 223}

The paper states (line~223):
\begin{quote}
  ``In theory, the $\lambda^\mathrm{ext}_\ell$ coefficients from both the
  firm-level and the sector-level should be the same, at least if each firm
  has equal weight in both cases.''
\end{quote}

\noindent\textbf{Assessment.}  This statement is \textbf{correct but
incomplete}.  Specifically:

\begin{enumerate}
  \item \textbf{The algebraic averaging is correct.}  Summing the
    firm-level equation over $\cFpre_{jme}$ and dividing by $N_c$ yields
    equation~\eqref{eq:sector} with the same population $\lambda^\mathrm{ext}_\ell$.
    No errors in the math.

  \item \textbf{``Equal weight'' must be enforced in \emph{both}
    regressions simultaneously.}  The condition ``each firm has equal weight''
    means:
    \begin{itemize}
      \item The firm-level regression is \emph{unweighted} (each firm-year
        observation gets weight~1).
      \item The sector-level regression uses analytic weights equal to the
        cell size $N_c = |\cFpre_{jme}|$.
    \end{itemize}
    If the sector-level regression is run unweighted (each cell gets weight~1),
    then firms in large cells are down-weighted relative to the firm-level
    regression, and $\hat\lambda^\mathrm{sector} \neq
    \hat\lambda^\mathrm{firm}$ in general.

  \item \textbf{The primary firm-level specification uses employment weights.}
    With employment weights at the firm level, equal-firm averaging at the
    sector level does not recover the same estimand.  To match, one would need
    employment-weighted sector averages and total-employment cell weights (see
    Section~3.2 above).  The paper wisely does not claim the primary
    (employment-weighted) firm-level and sector-level coefficients are the same;
    the claim is conditional on ``equal weight.''

  \item \textbf{Additional assumption: same sample.}  The aggregation identity
    requires that the sector-level regression uses \emph{exactly} the same firm
    set $\cFpre_{jme}$ as the firm-level regression.  If the sector regression
    drops cells that have no variation in $\FA$ (e.g., cells where all firms
    have zero affiliation), or if the firm panel has missing observations for
    some firms in some years, the samples diverge and the coefficients will
    differ even with correct weighting.

  \item \textbf{Additional assumption: same fixed-effect structure.}
    The firm-level regression absorbs $\gamma_f$ (firm FE) and $\alpha_{mt}$
    (municipality$\times$year FE).  The cell-averaged regression absorbs
    $\bar\gamma_{jme}$ (cell-level FE, which is a
    municipality$\times$sector$\times$election-cycle FE) and $\alpha_{mt}$.
    These are algebraically consistent \emph{only if} the cell-averaged
    regression is specified with the correct grouped FE.  In practice, the
    sector-level estimation in the paper uses $\gamma_{jm}$
    (municipality$\times$sector FE) and $\alpha_{jt}$ (sector$\times$year FE),
    which is a different FE structure from
    $\bar\gamma_{jme} + \alpha_{mt}$.  This means the sector-level estimator
    is \emph{not} a mechanical aggregation of the firm-level one; it is a
    complementary specification with its own identification.
\end{enumerate}

\section{Recommendation}

The claim on line~223 is mathematically sound for the population equation.
To make the paper's discussion precise, consider strengthening it to:

\begin{quote}
  ``In theory, if (i) the firm-level regression is unweighted, (ii)
  the sector-level regression weights each cell by $N^\mathrm{pre}_{jme}$,
  (iii) both regressions use the same sample of firms, and (iv) the
  sector-level fixed effects nest the averaged firm fixed effects,
  then the estimated $\lambda^\mathrm{ext}_\ell$ from the two regressions
  are numerically identical.  In practice, the sector-level specification
  uses a different fixed-effect structure
  ($\gamma_{jm} + \alpha_{jt}$ rather than $\bar\gamma_{jme} + \alpha_{mt}$),
  so the two estimators target related but distinct causal parameters.
  Comparing them is informative about the robustness of the political-lending
  channel across levels of aggregation.''
\end{quote}

\section{Relation to the Sector-Level Instrument $Z^\ell_{jmt}$}

A separate but related point: the sector-level instrument defined in
Section~2.3 of the paper is
\[
  Z^\ell_{jmt} = \sum_p w^\ell_{jmp,t}\,\Align^\ell_{mpt},
\]
where $w^\ell_{jmp,t}$ is the \emph{owner-count-weighted} average of
firm-level exposures:
\[
  w^\ell_{jmp,t}
  = \sum_{f\in\cF(j,m)} \frac{L_{f,t}}{N^\ell_{jm,t}}\,\omega^\ell_{fp,t}.
\]
This uses total-owner-count weighting, not equal-firm weighting.  The
``aggregation link'' paragraph (lines 330--347) confirms that $Z^\ell_{jmt}$
is the \emph{owner-count-weighted} average of $\FA^\ell_{fmt}$:
\[
  Z^\ell_{jmt}
  = \sum_{f\in\cFpre_{jme}} \frac{L_{f,t}}{\sum_{f'} L_{f',t}}\,\FA^\ell_{fmt}.
\]
This is distinct from the \emph{equal-firm average}
$\overline{\FA}^{\ell,\mathrm{pre}}_{jmt} = N_c^{-1}\sum_f \FA^\ell_{fmt}$
used in the extensive-margin aggregation.  The two objects would coincide only
if all firms in the cell have the same number of owners.

Therefore, the paper is internally consistent: the sector-level first stage
(Section~2.3) uses $Z^\ell_{jmt}$ as an instrument for $s_{jmt}$, while
Section~2.1's aggregation discussion uses
$\overline{\FA}^{\ell,\mathrm{pre}}_{jmt}$ as the sector counterpart of the
firm-level LPM.  These are different objects serving different purposes, and
the paper should not conflate them.

\end{document}
